\documentclass[aspectratio=169]{beamer}

\usetheme{colorful-dream}

\usepackage{minted}
\setminted{fontsize=\footnotesize, breaklines=true}

\title{Basics of Python for Data Science}
\subtitle{Lecture 3: Pandas, SciPy, and scikit-learn for Engineering Data}
\author{Lu Pa}
\date{\today}

% --- Title Slide Logos ---
\titlegraphic{
  \vspace{0.5cm}
  \begin{center}
    \includegraphics[height=1.8cm]{logo_university_placeholder.png}
    \hspace{1cm}
    \includegraphics[height=1.8cm]{logo_faculty_placeholder.png}
  \end{center}
}

\begin{document}

% ---------------------------------------------------------
\begin{frame}
  \titlepage
\end{frame}

% ---------------------------------------------------------
\section{PART I}
\begin{frame}[plain]
  \begin{center}
    \vspace{2cm}
    {\usebeamerfont{title}\color{white}\Huge PART I}
  \end{center}
\end{frame}

% ---------------------------------------------------------
\begin{frame}{What is Pandas?}
  \begin{itemize}
    \item High-level library for working with tabular data.
    \item Built on top of NumPy.
    \item Provides:
      \begin{itemize}
        \item DataFrame and Series structures
        \item fast filtering and grouping
        \item missing value handling
        \item merging and joining
      \end{itemize}
    \item Standard tool for data cleaning and reporting.
  \end{itemize}
\end{frame}

% ---------------------------------------------------------
\begin{frame}[fragile]{Loading Data with Pandas}
\begin{minted}{python}
import pandas as pd

df = pd.read_csv("../datasets/vehicle_speeds_trackA.csv")
df.info()
df.head()
\end{minted}
\end{frame}

% ---------------------------------------------------------
\begin{frame}{Cleaning Data}
  \begin{itemize}
    \item Remove impossible values:
      \begin{itemize}
        \item negative speed
        \item RPM spikes
      \end{itemize}
    \item Handle missing values:
      \begin{itemize}
        \item drop rows
        \item fill with interpolation
      \end{itemize}
    \item Convert units (km/h → m/s).
  \end{itemize}
\end{frame}

% ---------------------------------------------------------
\begin{frame}[fragile]{Filtering with Boolean Masks}
\begin{minted}{python}
high_speed = df[df["speed_kmh"] > 80]
braking = df[df["accel_mps2"] < -1.5]
cruise = df[(df["speed_kmh"] > 40) & (df["accel_mps2"].abs() < 0.2)]
\end{minted}
\end{frame}

% ---------------------------------------------------------
\begin{frame}[fragile]{Grouping and Aggregation}
\begin{minted}{python}
df["minute"] = (df["time_s"] // 60).astype(int)

summary = df.groupby("minute").agg({
    "speed_kmh": ["mean", "max"],
    "fuel_rate_gps": "mean",
    "engine_rpm": "mean"
})
\end{minted}
\end{frame}

% ---------------------------------------------------------
\begin{frame}{SciPy for Engineering Signals}
  \begin{itemize}
    \item SciPy extends NumPy with:
      \begin{itemize}
        \item signal processing
        \item optimization
        \item interpolation
        \item statistics
      \end{itemize}
    \item Today: smoothing and filtering.
  \end{itemize}
\end{frame}

% ---------------------------------------------------------
\begin{frame}[fragile]{Low-pass Filtering with SciPy}
\begin{minted}{python}
from scipy.signal import butter, filtfilt

b, a = butter(3, 0.05)  # 3rd order, cutoff 0.05
v_filtered = filtfilt(b, a, df["speed_kmh"])
\end{minted}
\end{frame}

% ---------------------------------------------------------
\begin{frame}{scikit-learn for Simple Modeling}
  \begin{itemize}
    \item Machine learning library for:
      \begin{itemize}
        \item regression
        \item classification
        \item clustering
      \end{itemize}
    \item Today: linear regression on telemetry.
  \end{itemize}
\end{frame}

% ---------------------------------------------------------
\begin{frame}[fragile]{Predicting Fuel Rate from Speed}
\begin{minted}{python}
from sklearn.linear_model import LinearRegression
import numpy as np

X = df[["speed_kmh"]].values
y = df["fuel_rate_gps"].values

model = LinearRegression()
model.fit(X, y)

print(model.coef_, model.intercept_)
\end{minted}
\end{frame}

% ---------------------------------------------------------
\begin{frame}[fragile]{Plotting Regression Fit}
\begin{minted}{python}
plt.scatter(df["speed_kmh"], df["fuel_rate_gps"], s=5)
plt.plot(df["speed_kmh"], model.predict(X), color="red")
plt.xlabel("Speed [km/h]")
plt.ylabel("Fuel Rate [g/s]")
plt.grid(True)
plt.show()
\end{minted}
\end{frame}

% ---------------------------------------------------------
\section{PART II}
\begin{frame}[plain]
  \begin{center}
    \vspace{2cm}
    {\usebeamerfont{title}\color{white}\Huge PART II}
  \end{center}
\end{frame}

% ---------------------------------------------------------
\begin{frame}{Engineering Workflow Tasks}
  \begin{itemize}
    \item Clean Track A and Track B datasets.
    \item Identify acceleration, braking, and cruising segments.
    \item Apply low-pass filtering to speed.
    \item Build a regression model predicting fuel rate.
    \item Compare Track A vs Track B driver behavior.
  \end{itemize}
\end{frame}

% ---------------------------------------------------------
\end{document}
